% The target is fixed-line-bundle Q-formal nef approximation, not an enlarged CLD cone.
\begin{theorem}
\label{thm:p081-irrational-tate-period-obstruction}
Let $K$ be an algebraically closed complete non-Archimedean field with
nontrivial real-valued absolute value, put $v=-\log|\cdot|$, and choose
$q,u\in K^\times$ such that
\[
  \ell:=v(q)>0,\qquad 0<a:=v(u)<\ell,
  \qquad a/\ell\notin\mathbf Q.
\]
Let $E_q/K$ be the Tate elliptic curve, let
$\pi:\mathbf G_{m,K}^{\mathrm{an}}\to E_q^{\mathrm{an}}$
be its uniformization with kernel $q^{\mathbf Z}$, let
$O=\pi(1)$, $P=\pi(u)$, and set
$L=\mathcal O_{E_q}(P-O)$.
Then the following assertions hold.
\begin{enumerate}
\item
$E_q$ is smooth, projective and geometrically connected, and $\deg L=0$.
For every integer $m>0$, there is no nef formal model of
$L^{\otimes m}$ on any proper admissible formal model of
$E_q^{\mathrm{an}}$. Here nef means nonnegative degree on every
irreducible component of the special fibre.
Consequently $L$ admits neither a semipositive
$\mathbf Q$-formal metric nor a numerically trivial
$\mathbf Q$-formal metric.
\item
After choosing a rigidification at the origin, the canonical cubist
metric $\|\cdot\|_{\mathrm{cub}}$ on $L$ is continuous and locally
CLD-psh, and its CLD curvature current satisfies
\[
 c_1(L,\|\cdot\|_{\mathrm{cub}})=0.
\]
In fact its local weights in the Tate uniformization are real affine
functions of $-\log|z|$, up to logarithms of analytic units, and are
therefore locally CLD-smooth and harmonic.
\item
The metric $\|\cdot\|_{\mathrm{cub}}$ is not a uniform limit of
semipositive $\mathbf Q$-formal metrics on this fixed line bundle
$L$. Thus local CLD-pshness does not imply membership in the
fixed-$L$ uniform closure of nef $\mathbf Q$-formal metrics over
arbitrary complete nontrivially valued fields, even when the field is
algebraically closed and the curve is smooth and projective.
\end{enumerate}
The hypotheses occur, for example, by choosing an irrational
$\alpha\in(0,1)$, equipping $\mathbf Q_p(T)$ with the weighted
Gauss absolute value extending $|p|=p^{-1}$ and satisfying
$|T|=p^{-\alpha}$, and taking the completed algebraic closure of its
completion. One may then take $q=p$, $u=T$.
\end{theorem}
